\hypertarget{concurrency}{%
\chapter{CONCURRENCY}\label{concurrency}}

\hypertarget{concurrency-multithreading}{%
\chapter{CONCURRENCY \&
MULTITHREADING}\label{concurrency-multithreading}}

C++11 brought a standard threading model.

\hypertarget{threads-stdthread}{%
\section{\texorpdfstring{1. Threads
(\texttt{std::thread})}{1. Threads (std::thread)}}\label{threads-stdthread}}

\begin{lstlisting}[language={C++}]
#include <thread>

void task() { /*...*/ }

int main() {
    std::thread t(task);
    t.join(); // Wait for finish
    // t.detach(); // Or let it run freely
}
\end{lstlisting}

\hypertarget{mutexes-locking}{%
\section{2. Mutexes \& Locking}\label{mutexes-locking}}

Avoid data races.

\begin{itemize}
\tightlist
\item
  \passthrough{\lstinline!std::mutex!}: Basic lock.
\item
  \passthrough{\lstinline!std::lock\_guard!}: RAII wrapper (locks on
  construction, unlocks on destruction).
\item
  \passthrough{\lstinline!std::unique\_lock!}: Flexible RAII wrapper
  (can unlock manually).
\end{itemize}

\begin{lstlisting}[language={C++}]
std::mutex mtx;
void safe() {
    std::lock_guard<std::mutex> lock(mtx);
    // critical section
}
\end{lstlisting}

\hypertarget{condition-variables}{%
\section{3. Condition Variables}\label{condition-variables}}

Wait for a condition to be true.

\begin{lstlisting}[language={C++}]
std::condition_variable cv;
std::mutex mtx;
bool ready = false;

// Waiter
std::unique_lock<std::mutex> lk(mtx);
cv.wait(lk, []{ return ready; });

// Notifier
{
    std::lock_guard<std::mutex> lk(mtx);
    ready = true;
}
cv.notify_one();
\end{lstlisting}

\hypertarget{futures-promises}{%
\section{4. Futures \& Promises}\label{futures-promises}}

Asynchronous result retrieval.

\begin{itemize}
\tightlist
\item
  \passthrough{\lstinline!std::async!}: Runs a function asynchronously.
\item
  \passthrough{\lstinline!std::future!}: Holds the result.
\end{itemize}

\begin{lstlisting}[language={C++}]
auto f = std::async(std::launch::async, []{ return 42; });
int result = f.get(); // Blocks until ready
\end{lstlisting}

\hypertarget{atomics-stdatomic}{%
\section{\texorpdfstring{5. Atomics
(\texttt{std::atomic})}{5. Atomics (std::atomic)}}\label{atomics-stdatomic}}

Lock-free operations for basic types.

\begin{lstlisting}[language={C++}]
std::atomic<int> counter(0);
counter++; // Thread-safe increment
\end{lstlisting}

\begin{center}\rule{0.5\linewidth}{0.5pt}\end{center}

\hypertarget{professional-notes-concurrency-depth}{%
\subsection{Professional Notes: Concurrency
Depth}\label{professional-notes-concurrency-depth}}

\hypertarget{thread-lifecycle-and-exceptions}{%
\subsubsection{1. Thread Lifecycle and
Exceptions}\label{thread-lifecycle-and-exceptions}}

If a \passthrough{\lstinline!std::thread!} object is destroyed while it
is still ``joinable'' (not joined or detached),
\passthrough{\lstinline!std::terminate()!} is called. * \textbf{Safety}:
Always use a wrapper or ensure
\passthrough{\lstinline!join()!}/\passthrough{\lstinline!detach()!} is
called in all exit paths, including exception handlers. *
\textbf{\passthrough{\lstinline!std::jthread!} (C++20)}: Automatically
joins on destruction, solving this problem.

\hypertarget{advanced-locking-strategies}{%
\subsubsection{2. Advanced Locking
Strategies}\label{advanced-locking-strategies}}

\begin{itemize}
\tightlist
\item
  \textbf{\passthrough{\lstinline!std::scoped\_lock!} (C++17)}: Locks
  multiple mutexes simultaneously using a deadlock-avoidance algorithm
  (replaces \passthrough{\lstinline!std::lock!}).
\item
  \textbf{\passthrough{\lstinline!std::shared\_mutex!} (C++17)}: Allows
  multiple readers or one writer (Reader-Writer Lock).
\item
  \textbf{Lock Strategies}:

  \begin{itemize}
  \tightlist
  \item
    \passthrough{\lstinline!std::adopt\_lock!}: Assume the calling
    thread already owns the mutex.
  \item
    \passthrough{\lstinline!std::defer\_lock!}: Do not lock the mutex on
    construction.
  \item
    \passthrough{\lstinline!std::try\_to\_lock!}: Attempt to lock
    without blocking.
  \end{itemize}
\end{itemize}

\hypertarget{semaphores-c20}{%
\subsubsection{3. Semaphores (C++20)}\label{semaphores-c20}}

A semaphore is a synchronization primitive that maintains a counter. *
\textbf{\passthrough{\lstinline!std::counting\_semaphore<N>!}}: Allows
up to \(N\) concurrent accesses. *
\textbf{\passthrough{\lstinline!std::binary\_semaphore!}}: Alias for
\passthrough{\lstinline!counting\_semaphore<1>!}. * \textbf{Usage}:
Useful for limiting access to a pool of resources (e.g., database
connections).

\hypertarget{thread-local-storage-tls}{%
\subsubsection{4. Thread Local Storage
(TLS)}\label{thread-local-storage-tls}}

The \passthrough{\lstinline!thread\_local!} keyword ensures that each
thread has its own unique instance of a variable.

\begin{lstlisting}[language={C++}]
thread_local int thread_id = 0; // Each thread gets its own copy
\end{lstlisting}

\begin{center}\rule{0.5\linewidth}{0.5pt}\end{center}

\hypertarget{professional-notes-chapter-80-threading}{%
\chapter{Professional Notes: Chapter 80:
Threading}\label{professional-notes-chapter-80-threading}}

Section 80.1: Creating a std::thread Section 80.2: Passing a reference
to a thread Section 80.3: Using std::async instead of std::thread
Section 80.4: Basic Synchronization Section 80.5: Create a simple thread
pool Section 80.6: Ensuring a thread is always joined Section 80.7:
Operations on the current thread Section 80.8: Using Condition Variables
Section 80.9: Thread operations Section 80.10: Thread-local storage
Section 80.11: Reassigning thread objects

\hypertarget{professional-notes-chapter-85-mutexes}{%
\chapter{Professional Notes: Chapter 85:
Mutexes}\label{professional-notes-chapter-85-mutexes}}

Section 85.1: Mutex Types Section 85.2: std::lock Section 85.3:
std::unique\_lock, std::shared\_lock, std::lock\_guard Section 85.4:
Strategies for lock classes: std::try\_to\_lock, std::adopt\_lock,
std::defer\_lock Section 85.5: std::mutex Section 85.6:
std::scoped\_lock (C++ 17)

\hypertarget{professional-notes-chapter-55-stdatomics}{%
\chapter{Professional Notes: Chapter 55:
std::atomics}\label{professional-notes-chapter-55-stdatomics}}

Section 55.1: atomic types Each instantiation and full specialization of
the std::atomic template denes an atomic type. If one thread writes to
an atomic object while another thread reads from it, the behavior is
well-dened (see memory model for details on data races) In addition,
accesses to atomic objects may establish inter-thread synchronization
and order non-atomic memory accesses as specied by std::memory\_order.
std::atomic may be instantiated with any TriviallyCopyable type T.
std::atomic is neither copyable nor movable. The standard library
provides specializations of the std::atomic template for the following
types: 1. One full specialization for the type bool and its typedef name
is dened that is treated as a non-specialized std::atomic except that it
has standard layout, trivial default constructor, trivial destructors,
and supports aggregate initialization syntax: Typedef name Full
specialization std::atomic\_bool std::atomic 2)Full specializations and
typedefs for integral types, as follows: Typedef name Full
specialization std::atomic\_char std::atomic std::atomic\_char
std::atomic std::atomic\_schar std::atomic std::atomic\_uchar
std::atomic std::atomic\_short std::atomic std::atomic\_ushort
std::atomic std::atomic\_int std::atomic std::atomic\_uint std::atomic
std::atomic\_long std::atomic std::atomic\_ulong std::atomic
std::atomic\_llong std::atomic std::atomic\_ullong std::atomic
std::atomic\_char16\_t std::atomic std::atomic\_char32\_t std::atomic
std::atomic\_wchar\_t std::atomic std::atomic\_int8\_t std::atomic
std::atomic\_uint8\_t std::atomic std::atomic\_int16\_t std::atomic
std::atomic\_uint16\_t std::atomic std::atomic\_int32\_t std::atomic
std::atomic\_uint32\_t std::atomic std::atomic\_int64\_t std::atomic
std::atomic\_uint64\_t std::atomic std::atomic\_int\_least8\_t
std::atomic std::atomic\_uint\_least8\_t std::atomic
std::atomic\_int\_least16\_t std::atomic std::atomic\_uint\_least16\_t
std::atomic std::atomic\_int\_least32\_t std::atomic
std::atomic\_uint\_least32\_t std::atomic std::atomic\_int\_least64\_t
std::atomic std::atomic\_uint\_least64\_t std::atomic
std::atomic\_int\_fast8\_t std::atomic std::atomic\_uint\_fast8\_t
std::atomic std::atomic\_int\_fast16\_t std::atomic
std::atomic\_uint\_fast16\_t std::atomic std::atomic\_int\_fast32\_t
std::atomic std::atomic\_uint\_fast32\_t std::atomic
std::atomic\_int\_fast64\_t std::atomic std::atomic\_uint\_fast64\_t
std::atomic std::atomic\_intptr\_t std::atomic std::atomic\_uintptr\_t
std::atomic std::atomic\_size\_t std::atomic std::atomic\_ptrdiff\_t
std::atomic std::atomic\_intmax\_t std::atomic std::atomic\_uintmax\_t
std::atomic Simple example of using std::atomic\_int \#include //
std::cout \#include // std::atomic, std::memory\_order\_relaxed
\#include // std::thread std::atomic\_int foo (0); void set\_foo(int x)
\{ foo.store(x,std::memory\_order\_relaxed); // set value atomically \}
void print\_foo() \{ int x; do \{ x =
foo.load(std::memory\_order\_relaxed); // get value atomically \} while
(x==0); std::cout \textless\textless{} ``foo:'' \textless\textless{} x
\textless\textless{} `\n'; \} int main () \{ std::thread first
(print\_foo); std::thread second (set\_foo,10); first.join();
//second.join(); return 0; \} //output: foo: 10

\hypertarget{professional-notes-chapter-80-threading-1}{%
\chapter{Professional Notes: Chapter 80:
Threading}\label{professional-notes-chapter-80-threading-1}}

Parameter other Details Takes ownership of other, other doesn't own the
thread anymore func args Function to call in a separate thread Arguments
for func Section 80.1: Creating a std::thread In C++, threads are
created using the std::thread class. A thread is a separate ow of
execution; it is analogous to having a helper perform one task while you
simultaneously perform another. When all the code in the thread is
executed, it terminates. When creating a thread, you need to pass
something to be executed on it. A few things that you can pass to a
thread: Free functions Member functions Functor objects Lambda
expressions Free function example - executes a function on a separate
thread (Live Example): \#include \#include void foo(int a) \{ std::cout
\textless\textless{} a \textless\textless{} `\n'; \} int main() \{ //
Create and execute the thread std::thread thread(foo, 10); // foo is the
function to execute, 10 is the // argument to pass to it // Keep going;
the thread is executed separately // Wait for the thread to finish; we
stay here until it is done thread.join(); return 0; \} Member function
example - executes a member function on a separate thread (Live
Example): \#include \#include class Bar \{ public: void foo(int a) \{
std::cout \textless\textless{} a \textless\textless{} `\n'; \} \}; int
main() \{ Bar bar; // Create and execute the thread std::thread
thread(\&Bar::foo, \&bar, 10); // Pass 10 to member function // The
member function will be executed in a separate thread // Wait for the
thread to finish, this is a blocking operation thread.join(); return 0;
\} Functor object example (Live Example): \#include \#include class Bar
\{ public: void operator()(int a) \{ std::cout \textless\textless{} a
\textless\textless{} `\n'; \} \}; int main() \{ Bar bar; // Create and
execute the thread std::thread thread(bar, 10); // Pass 10 to functor
object // The functor object will be executed in a separate thread //
Wait for the thread to finish, this is a blocking operation
thread.join(); return 0; \} Lambda expression example (Live Example):
\#include \#include int main() \{ auto lambda = \href{int\%20a}{} \{
std::cout \textless\textless{} a \textless\textless{} `\n'; \}; //
Create and execute the thread std::thread thread(lambda, 10); // Pass 10
to the lambda expression // The lambda expression will be executed in a
separate thread // Wait for the thread to finish, this is a blocking
operation thread.join();  return 0; \} Section 80.2: Passing a
reference to a thread You cannot pass a reference (or const reference)
directly to a thread because std::thread will copy/move them. Instead,
use std::reference\_wrapper: void foo(int\& b) \{ b = 10; \} int a = 1;
std::thread thread\{ foo, std::ref(a) \}; //`a' is now really passed as
reference thread.join(); std::cout \textless\textless{} a
\textless\textless{} `\n'; //Outputs 10 void bar(const ComplexObject\&
co) \{ co.doCalculations(); \} ComplexObject object; std::thread
thread\{ bar, std::cref(object) \}; //`object' is passed as const\&
thread.join(); std::cout \textless\textless{} object.getResult()
\textless\textless{} `\n'; //Outputs the result Section 80.3: Using
std::async instead of std::thread std::async is also able to make
threads. Compared to std::thread it is considered less powerful but
easier to use when you just want to run a function asynchronously.
Asynchronously calling a function \#include \#include unsigned int
square(unsigned int i)\{ return i*i; \} int main() \{ auto f =
std::async(std::launch::async, square, 8); std::cout
\textless\textless{} ``square currently running\n''; //do something
while square is running std::cout \textless\textless{} ``result is''
\textless\textless{} f.get() \textless\textless{} `\n'; //getting the
result from square \} Common Pitfalls std::async returns a std::future
that holds the return value that will be calculated by the function.
When that future gets destroyed it waits until the thread completes,
making your code eectively single threaded. This is easily overlooked
when you don't need the return value: std::async(std::launch::async,
square, 5); //thread already completed at this point, because the
returning future got destroyed std::async works without a launch policy,
so std::async(square, 5); compiles. When you do that the system gets to
decide if it wants to create a thread or not. The idea was that the
system chooses to make a thread unless it is already running more
threads than it can run eciently. Unfortunately implementations commonly
just choose not to create a thread in that situation, ever, so you need
to override that behavior with std::launch::async which forces the
system to create a thread. Beware of race conditions. More on async on
Futures and Promises Section 80.4: Basic Synchronization Thread
synchronization can be accomplished using mutexes, among other
synchronization primitives. There are several mutex types provided by
the standard library, but the simplest is std::mutex. To lock a mutex,
you construct a lock on it. The simplest lock type is std::lock\_guard:
std::mutex m; void worker() \{ std::lock\_guard guard(m); // Acquires a
lock on the mutex // Synchronized code here \} // the mutex is
automatically released when guard goes out of scope With
std::lock\_guard the mutex is locked for the whole lifetime of the lock
object. In cases where you need to manually control the regions for
locking, use std::unique\_lock instead: std::mutex m; void worker() \{
// by default, constructing a unique\_lock from a mutex will lock the
mutex // by passing the std::defer\_lock as a second argument, we // can
construct the guard in an unlocked state instead and // manually lock
later. std::unique\_lock guard(m, std::defer\_lock); // the mutex is not
locked yet! guard.lock(); // critical section guard.unlock(); // mutex
is again released \} More Thread synchronization structures Section
80.5: Create a simple thread pool C++11 threading primitives are still
relatively low level. They can be used to write a higher level
construct, like a thread pool: Version C++14 struct tasks \{ // the
mutex, condition variable and deque form a single // thread-safe
triggered queue of tasks: std::mutex m; std::condition\_variable v; //
note that a packaged\_task can store a packaged\_task: 
std::deque\textless std::packaged\_task\textless void()\textgreater\textgreater{}
work; // this holds futures representing the worker threads being done:
std::vector\textless std::future\textgreater{} finished; // queue(
lambda ) will enqueue the lambda into the tasks for the threads // to
use. A future of the type the lambda returns is given to let you get //
the result out. template\textless class F, class
R=std::result\_of\_t\textless F\&()\textgreater\textgreater{}
std::future queue(F\&\& f) \{ // wrap the function object into a
packaged task, splitting // execution from the return value:
std::packaged\_task\textless R()\textgreater{} p(std::forward(f)); auto
r=p.get\_future(); // get the return value before we hand off the task
\{ std::unique\_lock l(m); work.emplace\_back(std::move(p)); // store
the task\textless R()\textgreater{} as a
task\textless void()\textgreater{} \} v.notify\_one(); // wake a thread
to work on the task return r; // return the future result of the task \}
// start N threads in the thread pool. void start(std::size\_t N=1)\{
for (std::size\_t i = 0; i \textless{} N; ++i) \{ // each thread is a
std::async running this-\textgreater thread\_task():
finished.push\_back( std::async( std::launch::async, {[}this{]}\{
thread\_task(); \} ) ); \} \} // abort() cancels all non-started tasks,
and tells every working thread // stop running, and waits for them to
finish up. void abort() \{ cancel\_pending(); finish(); \} //
cancel\_pending() merely cancels all non-started tasks: void
cancel\_pending() \{ std::unique\_lock l(m); work.clear(); \} // finish
enques a ``stop the thread'' message for every thread, then waits for
them: void finish() \{ \{ std::unique\_lock l(m);
for(auto\&\&unused:finished)\{ work.push\_back(\{\}); \} \}
v.notify\_all(); finished.clear(); \} \textasciitilde tasks() \{
finish(); \} private: // the work that a worker thread does: void
thread\_task() \{ while(true)\{ // pop a task off the queue:
std::packaged\_task\textless void()\textgreater{} f; \{ // usual
thread-safe queue code: std::unique\_lock l(m); if (work.empty())\{
v.wait(l,{[}\&{]}\{return !work.empty();\}); \} f =
std::move(work.front()); work.pop\_front(); \} // if the task is
invalid, it means we are asked to abort: if (!f.valid()) return; //
otherwise, run the task: f(); \} \} \}; tasks.queue( {[}{]}\{ return
``hello world''s; \} ) returns a std::future, which when the tasks
object gets around to running it is populated with hello world. You
create threads by running tasks.start(10) (which starts 10 threads). The
use of packaged\_task\textless void()\textgreater{} is merely because
there is no type-erased std::function equivalent that stores move-only
types. Writing a custom one of those would probably be faster than using
packaged\_task\textless void()\textgreater. Live example. Version =
C++11 In C++11, replace result\_of\_t with typename result\_of::type.
More on Mutexes. Section 80.6: Ensuring a thread is always joined When
the destructor for std::thread is invoked, a call to either join() or
detach() must have been made. If a thread has not been joined or
detached, then by default std::terminate will be called. Using RAII,
this is generally simple enough to accomplish: class thread\_joiner \{
public: thread\_joiner(std::thread t) : t\_(std::move(t)) \{ \}
\textasciitilde thread\_joiner() \{ if(t\_.joinable()) \{ t\_.join(); \}
 \} private: std::thread t\_; \} This is then used like so: void
perform\_work() \{ // Perform some work \} void t() \{ thread\_joiner
j\{std::thread(perform\_work)\}; // Do some other calculations while
thread is running \} // Thread is automatically joined here This also
provides exception safety; if we had created our thread normally and the
work done in t() performing other calculations had thrown an exception,
join() would never have been called on our thread and our process would
have been terminated. Section 80.7: Operations on the current thread
std::this\_thread is a namespace which has functions to do interesting
things on the current thread from function it is called from. Function
Description get\_id Returns the id of the thread sleep\_for Sleeps for a
specied amount of time sleep\_until Sleeps until a specic time yield
Reschedule running threads, giving other threads priority Getting the
current threads id using std::this\_thread::get\_id: void foo() \{
//Print this threads id std::cout \textless\textless{}
std::this\_thread::get\_id() \textless\textless{} `\n'; \} std::thread
thread\{ foo \}; thread.join(); //`threads' id has now been printed,
should be something like 12556 foo(); //The id of the main thread is
printed, should be something like 2420 Sleeping for 3 seconds using
std::this\_thread::sleep\_for: void foo() \{
std::this\_thread::sleep\_for(std::chrono::seconds(3)); \} std::thread
thread\{ foo \}; foo.join(); std::cout \textless\textless{} ``Waited for
3 seconds!\n''; Sleeping until 3 hours in the future using
std::this\_thread::sleep\_until: void foo() \{
std::this\_thread::sleep\_until(std::chrono::system\_clock::now() +
std::chrono::hours(3)); \} std::thread thread\{ foo \}; thread.join();
std::cout \textless\textless{} ``We are now located 3 hours after the
thread has been called\n''; Letting other threads take priority using
std::this\_thread::yield: void foo(int a) \{ for (int i = 0; i
\textless{} al ++i) std::this\_thread::yield(); //Now other threads take
priority, because this thread //isn't doing anything important std::cout
\textless\textless{} ``Hello World!\n''; \} std::thread thread\{ foo, 10
\}; thread.join(); Section 80.8: Using Condition Variables A condition
variable is a primitive used in conjunction with a mutex to orchestrate
communication between threads. While it is neither the exclusive or most
ecient way to accomplish this, it can be among the simplest to those
familiar with the pattern. One waits on a std::condition\_variable with
a std::unique\_lock. This allows the code to safely examine shared state
before deciding whether or not to proceed with acquisition. Below is a
producer-consumer sketch that uses std::thread,
std::condition\_variable, std::mutex, and a few others to make things
interesting. \#include \#include \#include \#include \#include \#include
\#include int main() \{ std::condition\_variable cond; std::mutex mtx; 
std::queue intq; bool stopped = false; std::thread producer\{\href{}{\&}
\{ // Prepare a random number generator. // Our producer will simply
push random numbers to intq. // std::default\_random\_engine gen\{\};
std::uniform\_int\_distribution dist\{\}; std::size\_t count = 4006;\\
while(count--) \{\\
// Always lock before changing // state guarded by a mutex and //
condition\_variable (a.k.a. ``condvar''). std::lock\_guard L\{mtx\}; //
Push a random int into the queue intq.push(dist(gen)); // Tell the
consumer it has an int cond.notify\_one(); \} // All done. // Acquire
the lock, set the stopped flag, // then inform the consumer.
std::lock\_guard L\{mtx\}; std::cout \textless\textless{} ``Producer is
done!'' \textless\textless{} std::endl; stopped = true;
cond.notify\_one(); \}\}; std::thread consumer\{\href{}{\&} \{ do\{
std::unique\_lock L\{mtx\}; cond.wait(L,\href{}{\&} \{ // Acquire the
lock only if // we've stopped or the queue // isn't empty return stopped
\textbar\textbar{} ! intq.empty(); \}); // We own the mutex here; pop
the queue // until it empties out. while( ! intq.empty()) \{ const auto
val = intq.front(); intq.pop(); std::cout \textless\textless{}
``Consumer popped:'' \textless\textless{} val \textless\textless{}
std::endl; \} if(stopped)\{ // producer has signaled a stop  std::cout
\textless\textless{} ``Consumer is done!'' \textless\textless{}
std::endl; break; \} \}while(true); \}\}; consumer.join();
producer.join(); std::cout \textless\textless{} ``Example Completed!''
\textless\textless{} std::endl; return 0; \} Section 80.9: Thread
operations When you start a thread, it will execute until it is nished.
Often, at some point, you need to (possibly - the thread may already be
done) wait for the thread to nish, because you want to use the result
for example. int n; std::thread thread\{ calculateSomething, std::ref(n)
\}; //Doing some other stuff //We need `n' now! //Wait for the thread to
finish - if it is not already done thread.join(); //Now `n' has the
result of the calculation done in the separate thread std::cout
\textless\textless{} n \textless\textless{} `\n'; You can also detach
the thread, letting it execute freely: std::thread thread\{ doSomething
\}; //Detaching the thread, we don't need it anymore (for whatever
reason) thread.detach(); //The thread will terminate when it is done, or
when the main thread returns Section 80.10: Thread-local storage
Thread-local storage can be created using the thread\_local keyword. A
variable declared with the thread\_local specier is said to have thread
storage duration. Each thread in a program has its own copy of each
thread-local variable. A thread-local variable with function (local)
scope will be initialized the rst time control passes through its
denition. Such a variable is implicitly static, unless declared extern.
A thread-local variable with namespace or class (non-local) scope will
be initialized as part of thread startup. Thread-local variables are
destroyed upon thread termination. A member of a class can only be
thread-local if it is static. There will therefore be one copy of that
variable per thread, rather than one copy per (thread, instance) pair.
Example: void debug\_counter() \{ thread\_local int count = 0;
Logger::log(``This function has been called \%d times by this thread'',
++count); \} Section 80.11: Reassigning thread objects We can create
empty thread objects and assign work to them later. If we assign a
thread object to another active, joinable thread, std::terminate will
automatically be called before the thread is replaced. \#include void
foo() \{ std::this\_thread::sleep\_for(std::chrono::seconds(3)); \}
//create 100 thread objects that do nothing std::thread
executors{[}100{]}; // Some code // I want to create some threads now
for (int i = 0;i \textless{} 100;i++) \{ // If this object doesn't have
a thread assigned if (!executors{[}i{]}.joinable()) executors{[}i{]} =
std::thread(foo); \}

\hypertarget{professional-notes-chapter-85-mutexes-1}{%
\chapter{Professional Notes: Chapter 85:
Mutexes}\label{professional-notes-chapter-85-mutexes-1}}

Section 85.1: Mutex Types C++1x oers a selection of mutex classes:
std::mutex - oers simple locking functionality. std::timed\_mutex - oers
try\_to\_lock functionality std::recursive\_mutex - allows recursive
locking by the same thread. std::shared\_mutex,
std::shared\_timed\_mutex - oers shared and unique lock functionality.
Section 85.2: std::lock std::lock uses deadlock avoidance algorithms to
lock one or more mutexes. If an exception is thrown during a call to
lock multiple objects, std::lock unlocks the successfully locked objects
before re-throwing the exception. std::lock(\_mutex1, \_mutex2); Section
85.3: std::unique\_lock, std::shared\_lock, std::lock\_guard Used for
the RAII style acquiring of try locks, timed try locks and recursive
locks. std::unique\_lock allows for exclusive ownership of mutexes.
std::shared\_lock allows for shared ownership of mutexes. Several
threads can hold std::shared\_locks on a std::shared\_mutex. Available
from C++ 14. std::lock\_guard is a lightweight alternative to
std::unique\_lock and std::shared\_lock. \#include \#include \#include
\#include \#include \#include class PhoneBook \{ public: std::string
getPhoneNo( const std::string \& name ) \{ std::shared\_lock
l(\_protect); auto it = \_phonebook.find( name ); if ( it !=
\_phonebook.end() ) return (*it).second; return ""; \} void addPhoneNo (
const std::string \& name, const std::string \& phone ) \{
std::unique\_lock l(\_protect); \_phonebook{[}name{]} = phone; \}
std::shared\_timed\_mutex \_protect;
std::unordered\_map\textless std::string,std::string\textgreater{}
\_phonebook; \}; Section 85.4: Strategies for lock classes:
std::try\_to\_lock, std::adopt\_lock, std::defer\_lock When creating a
std::unique\_lock, there are three dierent locking strategies to choose
from: std::try\_to\_lock, std::defer\_lock and std::adopt\_lock 1.
std::try\_to\_lock allows for trying a lock without blocking: \{
std::atomic\_int temp \{0\}; std::mutex \_mutex; std::thread t(
\href{}{\&}\{ while( temp!= -1)\{
std::this\_thread::sleep\_for(std::chrono::seconds(5));
std::unique\_lock lock( \_mutex, std::try\_to\_lock);
if(lock.owns\_lock())\{ //do something temp=0; \} \} \}); while ( true )
\{ std::this\_thread::sleep\_for(std::chrono::seconds(1));
std::unique\_lock lock( \_mutex, std::try\_to\_lock);
if(lock.owns\_lock())\{ if (temp \textless{} INT\_MAX)\{ ++temp; \}
std::cout \textless\textless{} temp \textless\textless{} std::endl; \}
\} \} 2. std::defer\_lock allows for creating a lock structure without
acquiring the lock. When locking more than one mutex, there is a window
of opportunity for a deadlock if two function callers try to acquire the
locks at the same time: \{ std::unique\_lock lock1(\_mutex1,
std::defer\_lock); std::unique\_lock lock2(\_mutex2, std::defer\_lock);
lock1.lock() lock2.lock(); // deadlock here std::cout
\textless\textless{} "Locked! \textless\textless{} std::endl; //\ldots{}
\} With the following code, whatever happens in the function, the locks
are acquired and released in appropriate order: \{ std::unique\_lock
lock1(\_mutex1, std::defer\_lock); std::unique\_lock lock2(\_mutex2,
std::defer\_lock);  std::lock(lock1,lock2); // no deadlock possible
std::cout \textless\textless{} "Locked! \textless\textless{} std::endl;
//\ldots{} \} 3. std::adopt\_lock does not attempt to lock a second time
if the calling thread currently owns the lock. \{ std::unique\_lock
lock1(\_mutex1, std::adopt\_lock); std::unique\_lock lock2(\_mutex2,
std::adopt\_lock); std::cout \textless\textless{} "Locked!
\textless\textless{} std::endl; //\ldots{} \} Something to keep in mind
is that std::adopt\_lock is not a substitute for recursive mutex usage.
When the lock goes out of scope the mutex is released. Section 85.5:
std::mutex std::mutex is a simple, non-recursive synchronization
structure that is used to protect data which is accessed by multiple
threads. std::atomic\_int temp\{0\}; std::mutex \_mutex; std::thread t(
\href{}{\&}\{ while( temp!= -1)\{
std::this\_thread::sleep\_for(std::chrono::seconds(5));
std::unique\_lock lock( \_mutex); temp=0; \} \}); while ( true ) \{
std::this\_thread::sleep\_for(std::chrono::milliseconds(1));
std::unique\_lock lock( \_mutex, std::try\_to\_lock); if ( temp
\textless{} INT\_MAX ) temp++; cout \textless\textless{} temp
\textless\textless{} endl; \} Section 85.6: std::scoped\_lock (C++ 17)
std::scoped\_lock provides RAII style semantics for owning one more
mutexes, combined with the lock avoidance algorithms used by std::lock.
When std::scoped\_lock is destroyed, mutexes are released in the reverse
order from which they where acquired. \{ std::scoped\_lock
lock\{\_mutex1,\_mutex2\}; //do something \}

\begin{center}\rule{0.5\linewidth}{0.5pt}\end{center}

\hypertarget{chapter-16-concurrency-with-openmp}{%
\chapter{Chapter 16: Concurrency with
OpenMP}\label{chapter-16-concurrency-with-openmp}}

OpenMP (Open Multi-Processing) is an API that supports multi-platform
shared-memory multiprocessing programming in C, C++, and Fortran. It is
widely used in high-performance computing (HPC) for parallelizing loops
and sections of code with simple directives.

\hypertarget{getting-started-with-openmp}{%
\section{16.1 Getting Started with
OpenMP}\label{getting-started-with-openmp}}

OpenMP uses compiler directives (\passthrough{\lstinline!\#pragma omp!})
to parallelize code.

\hypertarget{parallel-regions}{%
\subsection{1. Parallel Regions}\label{parallel-regions}}

The most basic directive is
\passthrough{\lstinline!\#pragma omp parallel!}. It creates a team of
threads to execute the following block.

\begin{lstlisting}[language={C++}]
#include <iostream>
#include <omp.h>

int main() {
    #pragma omp parallel
    {
        int id = omp_get_thread_num();
        std::cout << "Hello from thread " << id << std::endl;
    }
    return 0;
}
\end{lstlisting}

\hypertarget{parallelizing-loops}{%
\subsection{2. Parallelizing Loops}\label{parallelizing-loops}}

OpenMP excels at parallelizing independent iterations of a loop.

\begin{lstlisting}[language={C++}]
#pragma omp parallel for
for (int i = 0; i < 1000; i++) {
    results[i] = compute(i);
}
\end{lstlisting}

\begin{center}\rule{0.5\linewidth}{0.5pt}\end{center}

\hypertarget{data-sharing-attributes}{%
\section{16.2 Data Sharing Attributes}\label{data-sharing-attributes}}

\begin{itemize}
\tightlist
\item
  \textbf{\passthrough{\lstinline!shared!}}: Variables are accessible by
  all threads.
\item
  \textbf{\passthrough{\lstinline!private!}}: Each thread has its own
  local copy of the variable.
\item
  \textbf{\passthrough{\lstinline!reduction!}}: Combines private copies
  into a single shared variable (e.g., sum, product).
\end{itemize}

\begin{lstlisting}[language={C++}]
double total = 0;
#pragma omp parallel for reduction(+:total)
for (int i = 0; i < 100; i++) {
    total += data[i];
}
\end{lstlisting}

\begin{longtable}[]{@{}
  >{\raggedright\arraybackslash}p{(\columnwidth - 0\tabcolsep) * \real{0.06}}@{}}
\toprule
\endhead
\#\#\# Professional Notes: OpenMP Performance \\
\#\#\#\# 1. Scheduling Strategies OpenMP provides different ways to
distribute loop iterations: * \passthrough{\lstinline!static!}:
Fixed-size chunks assigned at compile time (Low overhead). *
\passthrough{\lstinline!dynamic!}: Chunks assigned at runtime as threads
become free (Better for unbalanced workloads). *
\passthrough{\lstinline!guided!}: Chunks start large and shrink over
time to reduce tail-latency. \\
\#\#\#\# 2. False Sharing and Padding \textbf{Godhood Warning}: Avoid
``False Sharing,'' where multiple threads write to different variables
that happen to be on the same CPU cache line. This causes the cache line
to be repeatedly invalidated across cores, drastically reducing
performance. * \textbf{Fix}: Pad your data structures or ensure threads
work on data that is spaced apart in memory. \\
\#\#\#\# 3. Thread Affinity Use environment variables like
\passthrough{\lstinline!OMP\_PROC\_BIND=true!} to bind threads to
specific physical CPU cores, improving cache hits by preventing threads
from migrating between cores. \\
\bottomrule
\end{longtable}

\begin{center}\rule{0.5\linewidth}{0.5pt}\end{center}

\hypertarget{volume-03-refinement-generics-c14}{%
\chapter{VOLUME 03 REFINEMENT GENERICS
C14}\label{volume-03-refinement-generics-c14}}
